% Closed-loop wheel speed control block diagram.
%
% Forward path: target -> summing junction -> PI controller -> PWM/H-bridge -> motor.
% Feedback path: motor -> encoder -> speed calculation -> back into the summing
% junction with a minus sign. Pure vector TikZ (no opacity/shadings) so it survives
% the latex -> dvipdf build.
\begin{tikzpicture}[>=Latex, line join=round, font=\footnotesize,
    blk/.style={draw, line width=0.6pt, minimum height=1.0cm, minimum width=1.9cm,
                align=center},
    sum/.style={draw, circle, line width=0.6pt, minimum size=0.55cm, inner sep=0pt},
    sig/.style={->, line width=0.6pt}]

  % forward path
  \node[sum]                          (sum)   {};
  \node[blk, right=1.3cm of sum]      (pi)    {PI\\controller};
  \node[blk, right=1.3cm of pi]       (pwm)   {PWM +\\H-bridge};
  \node[blk, right=1.3cm of pwm]      (motor) {Motor\\+ wheel};

  % feedback path
  \node[blk, below=1.3cm of motor]    (enc)   {Encoder};
  \node[blk, below=1.3cm of pi]       (calc)  {Speed\\calculation};

  % input
  \coordinate[left=1.6cm of sum] (in);
  \draw[sig] (in) -- node[above] {$\dot{C}^*\,T_s$} node[below] {target} (sum);

  % forward arrows
  \draw[sig] (sum)   -- node[above] {$e[k]$}  node[below] {error} (pi);
  \draw[sig] (pi)    -- node[above] {$u[k]$}  node[below] {duty}  (pwm);
  \draw[sig] (pwm)   -- node[above] {voltage} (motor);

  % output tap
  \coordinate[right=0.9cm of motor] (out);
  \draw[sig] (motor) -- (out) node[right] {wheel speed};
  \draw[sig] ($(motor.east)!0.5!(out)$) |- (enc);

  % feedback arrows
  \draw[sig] (enc)  -- node[above] {pulses} (calc);
  \draw[sig] (calc) -| node[pos=0.75, left, align=right] {$\Delta C[k]$\\measured} (sum);

  % signs on the summing junction
  \node[above left=-0.05cm and 0.0cm of sum, font=\scriptsize] {$+$};
  \node[below left=0.0cm and -0.05cm of sum, font=\scriptsize] {$-$};
\end{tikzpicture}
